\mychapter{Desenho da Avaliação Experimental}{cap:experimental}\lhead{AVALIAÇÃO EXPERIMENTAL}

Este capítulo instancia a metodologia AIPE nas aplicações Bahia e M5. Dados, custos, prazos, políticas, modelos e janelas descritos a seguir são escolhas experimentais; não definem requisitos universais da metodologia. O protocolo separa o ajuste dos controladores, a construção do seletor e a avaliação das recomendações. A comparação entre abordagens do PSE examina a seleção contextual; os oráculos delimitam falhas de seleção, insuficiência do portfólio e inviabilidade física. As duas aplicações fornecem evidência sobre condições de uso, sem validar toda configuração admitida pela metodologia.

\section{Sistema de estoque e objetivo}
\label{sec:instanciacao_ambiente}\label{sec:ambiente_simulacao}
O ambiente simula vendas perdidas com prazo exato de dois ciclos. A política decide antes do recebimento; entregas chegam antes da demanda. Faltas não são acumuladas como pedidos pendentes. Quantidades não negativas podem ser fracionárias. Não há restrição conjunta de transporte, transbordo ou disputa entre lojas por estoque.

O custo da Equação~\ref{eq:obj_inv_control} soma manutenção do saldo após atendimento, falta, emissão e quantidade pedida. Os valores são $h=1$, $c_s=5$, $c_o^{\mathrm{fix}}=50$ e $c_o^{\mathrm{var}}=0{,}5$, sem crédito terminal. São unidades de cenário, sem calibração financeira. Pedidos em trânsito não atendem demanda antes da chegada.

\subsection{Métricas e definição do serviço}
\label{sec:metricas_avaliacao}
Registram-se custo total de inventário (CTI), atendimento (NS), taxa de ruptura (TR), frequência de pedidos (FP) e amplificação de variância (BE):
\[
 \mathrm{NS}=\frac{\sum_t(D_t-S_t)}{\sum_tD_t},\qquad
 \mathrm{TR}=\frac{\#\{t:S_t>0\}}T,\qquad
 \mathrm{FP}=\frac{\#\{t:Q_t>0\}}T.
\]
Aqui, $S_t$ é falta e BE é $\mathrm{Var}(Q)/\mathrm{Var}(D)$, com convenção registrada para variância nula. NS é fração de unidades atendidas, não probabilidade de um ciclo sem falta. Demanda total zero mantém NS indefinido; o custo continua na análise.

O limiar é 0,70. A média por série, a razão entre unidades totais atendidas e demandadas e a garantia de 0,70 em cada série não são equivalentes. Por isso, reportamos custo, média de NS, atendimento ponderado quando disponível e contagem de violações. Custo baixo com pouco atendimento não caracteriza eficiência operacional.

\section{Dados, preparo e cobertura}
\label{sec:pipeline_experimental}\label{sec:ba2026_methods}
A avaliação proprietária utiliza 145 séries da Bahia, em 115 unidades comerciais e nove produtos, com 38 ciclos de 202301 a 202504. As unidades correspondem aos revendedores identificados na fonte e são denominadas lojas nesta avaliação. Registros complementares são somados por armazém, loja, produto e ciclo; células sem agregado registrado são preenchidas com zero, sob hipótese explícita de cobertura da fonte. A coorte mantém identificadores históricos. Em cada origem, exige-se pelo menos três ciclos positivos no treino. Esse requisito é passado-dependente, mas não torna prospectiva a seleção inicial dos identificadores.

O M5 utiliza a coorte histórica local de 2.408 identificadores loja--produto da categoria de alimentos (\texttt{FOODS}), nas dez lojas de três estados. Usam-se ciclos completos de 21 dias, sem corte de picos pelo intervalo interquartil (\textit{interquartile range}, IQR). A exigência de três ciclos positivos no prefixo admite 885 séries na origem 50, 2.033 na origem 64 e 2.408 na origem 78. O recorte não representa todas as séries FOODS. Um piloto de 30 séries orientou o desenvolvimento; sua ampliação não estabelece confirmação independente.

\begin{table}[htbp]\centering\small
\caption{Matrizes completas incluídas na avaliação. Sementes e origens são medidas repetidas.}\label{tab:temporal_coverage}
\begin{tabular}{@{}lrr@{}}\toprule
Dimensão & Bahia & M5\\\midrule
Séries / lojas & 145 / 115 & 2.408 / 10\\
Origens do treino & 17 e 23 & 50, 64 e 78\\
Validação / teste & 7 / 7 ciclos & 7 / 7 ciclos\\
Casos série--origem & 290 & 5.326\\
Ajustes de políticas & 4.350 & 79.890\\
Testes com demanda positiva & 283 & 5.138\\\bottomrule
\end{tabular}\end{table}

A auditoria da coorte canônica Bahia identificou 2.697 de 5.510 células preenchidas com zero por ausência de agregado (48,9\%). Essa contagem descreve o painel auditado, sem estimar demanda perdida nem caracterizar a base M5. Vendas podem omitir demanda censurada por ruptura. Reindexar ausência de venda como zero pressupõe cobertura da fonte, não prova ausência de procura. Os dados brasileiros são restritos; M5 é público. O recorte local é inspirado no contexto de \citet{Zabraoui2025GaDQN}, sem identidade comprovada de filtros, amostra ou implementação.

\section{Características e perfis}
\label{sec:carac_series}
O preditivo recebe 17 características: ADI, CV$^2$, intermitência, proporção de zeros, maior e média das sequências de zeros, burstiness, assimetria, curtose, entropia, tendência, autocorrelação, média, desvio-padrão, CV, extensão e número de ciclos positivos. Todas usam somente o treino; perfil e metadados não alteram essa restrição temporal.

A característica CV$^2$ inclui zeros, diferindo da taxonomia original baseada em tamanhos positivos (Seção~\ref{sec:sb_taxa}). Burstiness é $(s_\tau-\bar\tau)/(s_\tau+\bar\tau+10^{-9})$ nos intervalos entre ocorrências positivas, recebendo zero com menos de dois intervalos \citep{GohBarabasi2008Burstiness}. Tendência é a inclinação por mínimos quadrados ordinários (\textit{ordinary least squares}, OLS) dividida pela média regularizada; entropia usa dez classes de mesma largura. Autocorrelação na defasagem 17 recebe zero com até 18 ciclos ou variância desprezível. Zero nessa condição significa estimativa indisponível.

\begin{table}[htbp]\centering\small
\caption{Perfis: primeira regra satisfeita determina a classe.}\label{tab:pod}
\begin{tabular}{@{}p{4.2cm}p{8.2cm}@{}}\toprule
Perfil & Condição\\\midrule
Sparse High Impact & ADI $\geq1{,}32$, CV$^2\geq0{,}49$, sequência de zeros $\geq5$\\
High Volatility Seasonal & Autocorrelação $\geq0{,}25$ ou burstiness $\geq0{,}30$\\
Unstable Trend & Tendência absoluta $\geq0{,}05$ ou entropia $\geq2$\\
Low Volume Stable & Média $\leq2$, CV$^2<0{,}49$\\
Fast Moving & ADI $<1{,}32$\\\bottomrule
\end{tabular}\end{table}

Os nomes são heurísticos; casos não cobertos recebem Sparse High Impact. Não se assume sazonalidade ou impacto econômico empiricamente comprovados. A coluna de política previamente associada aos nomes no código não gera os rótulos do experimento: custos e serviço na validação definem a escolha.

\section{Portfólio e calibração}
\label{sec:benchmark_massivo}\label{sec:portfolio_ampliado}
EOQCycle, sSCycle, CappedBaseStock, GA, PPO, GA--DQN e DQN compõem o portfólio. DQN é controle necessário do híbrido. Outros métodos preservados no projeto não são resultados adicionais desta matriz. A composição privilegia diversidade de regras; sua retirada do escopo não estabelece inferioridade universal.

\subsection{Referências e busca de parâmetros}
\label{sec:sinais_previsao}\label{sec:redundancia_portfolio}
Com média $\hat\mu$ e desvio populacional $\hat\sigma$ do treino, as escalas são
\begin{align}
 I_0&=\hat\mu L+1{,}28\hat\sigma\sqrt L,\label{eq:I0_metod}\\
 Q_{\mathrm{ref}}&=\max(1,\sqrt{2\hat\mu c_o^{\mathrm{fix}}/h}),\label{eq:qref_metod}\\
 s_{\mathrm{ref}}&=\hat\mu L+1{,}28\hat\sigma\sqrt L.\label{eq:sref_metod}
\end{align}
EOQCycle pede lote fixo ao atingir o ponto; sSCycle repõe até $s_{\mathrm{ref}}+Q_{\mathrm{ref}}$. Não se alega otimização exata de $(s,S)$. CappedBaseStock calibra alvo e teto em grade $20\times12$, motivada pela família de \citet{Xin2021CappedBaseStock}.

GA ajusta gatilho, lote e proteção adicional com 100 indivíduos e 50 gerações. Seu objetivo acrescenta penalidade pelo déficit de serviço à referência
\begin{equation}
 \boldsymbol\theta_i^*=\arg\min_{\boldsymbol\theta\in\Theta_i}\mathrm{CTI}(\pi_{\boldsymbol\theta},i)
 \quad\text{sujeito a }\mathrm{NS}(\pi_{\boldsymbol\theta},i)\geq0{,}70.
 \label{eq:obj_metaheuristicas}
\end{equation}
Penalização não garante viabilidade; o serviço observado deve ser conferido. A CappedBaseStock utiliza pesos zero para ruptura e amplificação dos pedidos nas duas bases. A harmonização do objetivo permite comparar as instanciações sob a mesma definição de calibração, preservando diferenças de demanda e escala.

\subsection{Reforço e híbrido}
DQN utiliza Double DQN dueling e PPO utiliza otimização proximal, com duas camadas de 64 unidades e 20 ações discretas. O estado reúne estoque, referência ingênua de demanda, pedidos em trânsito, média e dispersão dos últimos seis ciclos e posição no horizonte. Escalas vêm do treino; observações são passadas. Previsores especializados disponíveis no projeto não alimentam os resultados reportados.

O orçamento mínimo é de 20 mil transições, arredondado ao episódio completo. GA--DQN acrescenta a busca genética e mil demonstrações. Acima do gatilho genético, pede zero; abaixo, executa o máximo entre a ação da rede e metade do lote genético. Essa transformação é igual em demonstração, treino e aplicação. O esforço excede o do DQN; não se isola causalmente o benefício da combinação \citep{Zabraoui2025GaDQN}.

GA, PPO, DQN e GA--DQN usam sementes 11, 29 e 47; as referências determinísticas têm um ajuste. Primeiro promediamos as sementes dentro de série, origem e política. Elas representam variação do treinamento, não novas realizações de demanda.

\section{Separação temporal e comparação}
\label{sec:simulacao_benchmark}\label{sec:wf_protocol}
Em cada origem $o$, o desenho é
\begin{equation}
 \underbrace{1,\ldots,o}_{\text{treino}}\ ;\
 \underbrace{o+1,\ldots,o+7}_{\text{validação}}\ ;\
 \underbrace{o+8,\ldots,o+14}_{\text{teste}}.
 \label{eq:wf_splits}
\end{equation}
Na Bahia, os testes compartilham um ciclo; no M5, ocupam 58--64, 72--78 e 86--92, sem sobreposição. Não se combinam origens posteriores para recomendar em uma anterior.

A validação começa após seis ciclos de aquecimento com referência EOQ. A trajetória dessa referência na validação fornece estoque e pipeline comuns no teste. O estado usa apenas o passado e não continua o estoque próprio de cada candidata. Essa convenção isola comparações no mesmo ponto inicial, sem simular implantação contínua de cada política.

Uma loja inteira é excluída do conjunto que define o seletor, separadamente em cada origem. As outras lojas fornecem resultados de validação; a loja reservada fornece o teste. O controlador conserva o treino no histórico da própria série. Avalia-se transferência do seletor com histórico disponível, não generalização para lojas sem observações.

\section{Instanciação dos seletores}
\label{sec:selecao_contextual}
\subsection{Associação e rótulos}
Para um conjunto $A$ de casos de validação, a regra global escolhe
\begin{equation}
 \hat\pi_A=\arg\min_{\pi:\bar F_{\pi,A}\geq0{,}70}\bar C_{\pi,A}.
 \label{eq:policy_selector}
\end{equation}
O PSE por perfil aplica a regra a cada $A_g$. Sem viável, escolhe maior NS e menor custo; perfis sem suporte usam a global. O rótulo preditivo usa o mesmo critério na série individual. Serviço inteiramente indefinido leva a menor custo, com marcação de demanda zero. Empates restantes usam convenção determinística registrada na análise.

As restrições por série e por grupo diferem; a moda dos rótulos é um controle adicional com o mesmo alvo categórico. Todos os seletores são avaliados pelas consequências de suas escolhas. Acurácia é auxiliar: errar entre políticas quase empatadas difere de perder muito serviço.

\subsection{Aprendizado e recuo conservador}
O preditivo por classe utiliza XGBoost com 200 árvores, profundidade três, taxa 0,05, semente fixa e sem subamostragem. Imputação, codificação do perfil e classes são ajustadas dentro de cada fold. Classes raras não são removidas; uma classe única leva a classificador constante. A quantidade de exemplos de treinamento depende da origem e da loja excluída.

A ablação retira somente a entrada de perfil. O piloto anterior motivou alternativas de aprendizagem e intervenção; a avaliação consolidada aplica os seletores às matrizes completas e mantém a distinção entre desenvolvimento e confirmação independente \citep{CawleyTalbot2010}.

\subsection{Comparação por consequências}
A comparação usa a mesma matriz de cada base, sem repetir ajustes de controladores. Duas florestas ExtraTrees estimam custo e serviço das sete candidatas: 64 árvores, profundidade máxima quatro e mínimo de três observações por folha. O custo é aprendido como $\log(1+C)$ e transformado de volta. O serviço usa casos definidos, sem transformar demanda zero em atendimento completo. A recomendação direta escolhe menor custo previsto entre candidatas com serviço previsto de pelo menos 0,70; sem candidata viável, prioriza serviço e depois custo.

A variante com recuo usa a global como referência e exige serviço previsto de pelo menos 0,70, serviço previsto não inferior ao da referência e economia acima de uma margem. A margem pertence a $\{0;0{,}02;0{,}05;0{,}10;\infty\}$, sendo $\infty$ o recuo permanente. Ela minimiza custo na validação interna sujeita a serviço médio não inferior à referência. Cada origem mantém uma loja externa reservada e divide as demais em cinco grupos de lojas. Cada grupo interno recebe previsões de um modelo que exclui suas observações. Depois da escolha da margem, o modelo externo é ajustado em todas as lojas admitidas e recomenda para a loja reservada.

Um classificador ExtraTrees com a mesma regularização oferece um controle por rótulo vencedor; outra ablação remove o perfil da previsão de consequências sem recuo. As recomendações são persistidas antes da pontuação do teste. O piloto orientou o desenho; a análise consolidada não é apresentada como confirmação independente. O critério empírico não garante atendimento futuro, e a diferença entre objetivos individuais previstos e objetivos médios dos comparadores permanece uma limitação.

\subsection{Execução e compatibilidade das matrizes}
A execução separa planejamento, ajuste dos controladores e avaliação dos seletores. A preparação audita painéis, identidades, código e configurações antes de aproveitar checkpoints. A avaliação recusa matrizes incompletas e usa o mesmo pré-processamento para modelos por classe e por consequências, ajustado somente no conjunto admitido em cada partição. As duas bases incluem global, perfil, classe vencedora com e sem perfil, moda dos rótulos, floresta de classe de controle e previsão de consequências com e sem perfil e com recuo. Oráculos, intervenções e deteriorações completam a avaliação.

A Bahia reúne 4.350 ajustes e o M5, 79.890. Ajustes previamente compatíveis foram reaproveitados após auditoria, sem serem contados como novas observações independentes. A avaliação brasileira preserva sua relevância aplicada, mas continua sendo simulação com custos de cenário e possível censura das vendas. As bases fornecem contextos complementares, sem combinar custos entre elas.

\subsection{Aplicação e reavaliação}
\label{sec:aplicacao_aipe}
O uso pretendido entrega recomendação, parametrização e evidências para decisão do gerente. Na avaliação atual, a família escolhida permanece por sete ciclos de teste, com ações de reposição em cada período. Essa configuração não valida uma nova seleção a cada ciclo comercial. Uma avaliação sequencial deverá separar revisão do PSE, atualização do seletor e recalibração dos controladores, mantendo estado e pedidos em trânsito próprios de cada estratégia. Integração operacional e supervisão humana exigem validação adicional.

\section{Oráculos e decomposição das falhas}
O melhor controlador único no teste e a melhor candidata por caso são referências retrospectivas. Não fornecem rótulos para a escolha. O oráculo do portfólio usa as médias entre sementes e a regra de menor custo viável. Empates numéricos contam como acerto sem exigir identidade nominal.

Nas duas bases, o planejamento com informação perfeita conhece a demanda do horizonte. Sua programação linear inteira mista (\textit{mixed-integer linear programming}, MILP) mantém custos, prazo e estado inicial, com quantidades contínuas e variáveis binárias de emissão e saldo. Para falta $u_t$, chegada inicial $a_t$ e pedido $q_{t-L}$:
\[
 I_t-u_t=I_{t-1}+a_t+q_{t-L}-d_t,\quad
 I_tu_t=0,\quad \sum_tu_t\leq(1-\alpha)\sum_td_t.
\]
A complementaridade é linearizada com variáveis binárias e limites suficientes para o horizonte. Pedidos sem chegada antes do fim recebem zero, sem crédito terminal. Quando o alvo é fisicamente impossível nas primeiras entregas, maximiza-se atendimento antes de minimizar custo. Cada solução é conferida no simulador e recebe status de otimalidade.

Esse ótimo é determinístico e condicional, não o melhor controlador estocástico implantável. A eficiência $100C^*/C_{\mathrm{AIPE}}$ só é calculada sob o mesmo critério de serviço, com ambos os custos zero equivalendo a 100\%. Compara-se o MILP com cada realização de treinamento viável, pois cumprir o alvo na média das sementes não implica cumpri-lo em cada uma.

A decomposição distingue atendimento do alvo, erro de seleção com alternativa viável, ausência de candidata viável apesar de possibilidade física, limite inicial do sistema e demanda zero. Assim, uma falta não é atribuída automaticamente ao seletor.

\section{Incerteza e rastreabilidade}
\label{sec:protocolo_exp}\label{sec:reporting}\label{app:reproducibility}
Bootstrap pareado preserva produtos e origens de cada loja, com 5.000 sorteios em cada base. Ajustes e escolhas permanecem fixos; os intervalos são condicionais. Dependência entre lojas que compartilham produtos e a quantidade de somente dez lojas do M5 limitam a inferência, mesmo com milhares de casos série--origem.

A auditoria das entradas encontrou vendas observadas nas duas bases, mas custos, prazo e estados iniciais de cenário. Preços e receita disponíveis não calibram custos logísticos. A fotografia brasileira de estoque, de 21/05/2025, é posterior aos testes e usa vínculos múltiplos entre franquias e revendedores; cobertura de códigos não comprova alocação física ou saldo histórico. Ela permanece fora dos estados iniciais.

Manifestos, snapshots, sementes e bancos transacionais identificam cada matriz. As matrizes são:\par\begingroup\small\raggedright \nolinkurl{aipe_ba145_unified_2026-10-15} e
\nolinkurl{aipe_m5_2408_unified_2026-10-15}\par\endgroup
A pasta comum é \nolinkurl{experiment/data/08_reporting/research_v2}. As recomendações são gravadas antes da pontuação externa. Dados não integram o Git e possuem backup separado; código e fontes documentais estão versionados.

A rastreabilidade existente é parcial: configurações, sementes, indicadores e recomendações são preservados, mas faltam parâmetros finais de todas as políticas, redes reutilizáveis, trajetórias e custos desagregados por caso. Uma prévia gerencial da Bahia consolidada separa recomendação e alternativas de validação dos resultados posteriores, com ausências explícitas; não completa esses artefatos. Recuperar detalhes pode exigir replay ou novo ajuste, registrado separadamente.

Os estudos brasileiros anteriores, incluindo a qualificação, motivam a investigação. A evidência quantitativa deste capítulo utiliza somente as matrizes consolidadas descritas, sem misturar configurações. A expansão nacional está cancelada. O módulo opcional de atributos adicionais é examinado em estudo de desenvolvimento separado, reutilizando as matrizes de controladores. As variantes desse estudo compartilham floresta de 32 árvores e regra de intervenção fixa; não são contrastes isolados de atributos frente aos modelos de 64 árvores ou ao XGBoost da avaliação principal. A ativação depende da disponibilidade temporal documentada. Na Bahia, a correspondência do cadastro com o período é informada pelo autor, sem histórico de alterações por campo.
